\chapter{The Arakelov Volume Theorem on Unitary Flag Towers in Signature $(p, q)$}
\label{ch:13}

To establish the upper covolume bound on the dual automorphic row lattice $\mathcal{N}^\vee$ (\cref{ch:14}), we must bound from below the \emph{Arakelov mean height} of canonical automorphic vector bundles over unitary Shimura varieties of indefinite signature $(p, q)$ ($q \ge 2$, $p \ge 2q + 1$). Classically, arithmetic Hilbert--Samuel theory and Moriwaki's arithmetic bigness theorem~\cite{GilletSoule1992,BostGilletSoule1994,Moriwaki2009,Yuan2008,Zhang1995} relate the mean height of a Hermitian line bundle to its top arithmetic self-intersection number—but for high-rank automorphic vector bundles of well-separated highest weights on non-compact PEL Shimura varieties, boundary cusp contributions and wall crossings can easily spoil naive positivity.

This chapter proves the unconditional \emph{Arakelov Mean Height Lower Bound in Signature $(p, q)$} (\cref{thm:13:volume-pq}) whenever $p \ge 2q + 1$. We organize the geometry cleanly around three pillars:
\begin{enumerate}[label=(\roman*)]
\item a transparent \emph{Six-Cluster Classical Geometric Box} consolidating all $42$ published algebro-geometric, Arakelov-theoretic, and $p$-adic Hodge inputs (zero omissions, zero dead branches);
\item the \emph{Isotropic Counting Bootstrap} showing precisely why $p \ge 2q + 1$ ($\kappa = 2f(p+q-1) \ge 3a = 6fq$) guarantees strict dominance of the interior flag volume over all toroidal boundary cusps; and
\item a streamlined \emph{Four-Phase Abstract Flag-Bundle Tower Engine} (distilling the archive's eight-step recursion while preserving every explicit constant $K_{\mathrm{tot}} = 2C + 3C' + 5$, $C' = 4$, $C_u = 4$) instantiated on the PEL unitary Shimura root tower.
\end{enumerate}

\begin{remark}[Explicit Unsquared-Slope Convention Warning for \cref{ch:13}]
\label{rem:13:unsquared-warning}
Throughout \cref{ch:13} alone, all arithmetic intersection numbers $\widehat{\deg}$, component degrees $\widehat{\deg}_Z(u_M)$, and flag-tower mean heights $\widehat{\mu}_{\mathcal{T}}(B, c; u)$ follow the standard \emph{unsquared} Gillet--Soul\'e / Bost Arakelov normalization~\cite{GilletSoule1992,BostGilletSoule1994} ($\widehat{\mu}_{\mathcal{T}} = \frac{1}{2}\mu$ up to field-discriminant shift). The exact factor-of-$2$ conversion back to our monograph's squared-norm slope $\mu$ of \eqref{eq:00:squared-slope} is executed explicitly at four marked checkpoints in \cref{ch:14}.
\end{remark}

\section{The Six-Cluster Classical Geometric Box (42 Published Inputs)}
\label{sec:13:six-clusters}

In the formal archive, the geometric inputs to the volume theorem and covolume descent were split across $15$ top-level geometric axioms and $28$ Hodge-bundle fields ($42$ distinct published leaf theorems after removing internal proved deductions). To give the reader an unobstructed view of every classical theorem invoked in \cref{ch:13,ch:14}, we group all $42$ inputs into \emph{six thematic clusters} (the full $42$-item bijection table is audited in \cref{app:A}).

\newsimplification{Six-Cluster Thematic Consolidation of All 42 Published Geometric Inputs}

\begin{proposition}[Six-cluster classical geometric box]
\label{prop:13:six-clusters}
Over our CM field $E/E^+$ and its integral PEL unitary Shimura varieties of signature $(p, q)$ ($q \ge 2$, $p \ge 2q + 1$), the following six published geometric packages hold unconditionally:
\begin{enumerate}[label=\textup{\textbf{Cluster~\Roman*.}}]
\item \emph{(Resolution, principalization, and complex/formal GAGA — $5$ inputs.)}
  Characteristic-$0$ resolution of singularities and ideal principalization (\litref{Hironaka (1964)~\cite{Hironaka1964}}), Serre's GAGA comparison of algebraic and holomorphic global sections (\litref{Serre (1956)~\cite{Serre1956}}), and Grothendieck's formal algebraization theorem.

\item \emph{(Arakelov intersections, Moriwaki bigness, arithmetic Hilbert--Samuel, and cycle heights — $11$ inputs.)}
  Poincar\'e--Lelong inversion, arithmetic intersection calculus, Green-series expansions, Chow forms, cycle heights, arithmetic Kleiman positivity, the arithmetic Hilbert--Samuel theorem, and Moriwaki's arithmetic bigness criterion (\litref{Gillet--Soul\'e (1992)~\cite{GilletSoule1992}; Bost--Gillet--Soul\'e (1994)~\cite{BostGilletSoule1994}; Zhang (1995)~\cite{Zhang1995}; Yuan (2008)~\cite{Yuan2008}; Moriwaki (2009)~\cite{Moriwaki2009}}).

\item \emph{(Complex algebraic positivity, vanishing, $L^2$-extension, and point separation — $6$ inputs.)}
  Nef Riemann--Roch and numerical bigness, Kawamata--Viehweg vanishing (\litref{Kawamata (1982)~\cite{Kawamata1982}}), Ohsawa--Takegoshi $L^2$-extension for log-canonical metrics (\litref{Ohsawa--Takegoshi (1987)~\cite{OhsawaTakegoshi1987}}), Nadel vanishing and non-klt multiplier-ideal divisors (\litref{Nadel (1990)~\cite{Nadel1990}}), and Demailly / Angehrn--Siu effective point separation (\litref{Demailly (1993)~\cite{Demailly1993}; Angehrn--Siu (1995)~\cite{AngehrnSiu1995}}).

\item \emph{(Integral toroidal compactifications and semi-abelian geometry — $9$ inputs.)}
  Integral smooth projective toroidal compactifications of Siegel and PEL unitary Shimura varieties, semistable reduction, extension of semi-abelian homomorphisms, Mumford's good singular Chern--Weil metrics on automorphic bundles, seminegativity of semi-abelian Lie algebras, and principal-polarization section normalization (\litref{Mumford (1977)~\cite{Mumford1977}; Moret-Bailly (1985)~\cite{MoretBailly1985}; Faltings--Chai (1990)~\cite{FaltingsChai1990}; Lan (2013)~\cite{Lan2013}; SGA~7~\cite{SGA7}}).

\item \emph{($p$-adic Hodge theory, prismatic Dieudonn\'e theory, Fargues exactness, and deformation — $6$ inputs.)}
  Fargues's Harder--Narasimhan exactness / Frobenius prime-jump formula for finite flat group schemes (\litref{Fargues (2010)~\cite{Fargues2010}}), Zarhin's trick for principal polarizations (\litref{Zarhin (1974)~\cite{Zarhin1974}}), Grothendieck--Messing crystalline deformation theory (\litref{Messing (1972)~\cite{Messing1972}}), abelian crystalline--de~Rham comparison, and prismatic Dieudonn\'e theory (\litref{Ansch\"utz--Le~Bras (2023)~\cite{AnschutzLeBras2023}}).

\item \emph{(Class field theory, Hermitian forms, strong approximation, and modular representations — $5$ inputs.)}
  Ray class fields (\litref{Artin (1927)~\cite{Artin1927}}), strong approximation for simply connected special unitary groups $\SU(V)$ (\litref{Kneser (1965)~\cite{Kneser1965}; Platonov--Rapinchuk (1994)~\cite{PlatonovRapinchuk1994}}), Landherr's classification and the Hasse norm theorem (\litref{Landherr (1936)~\cite{Landherr1936}}), and induced simple modules in the lowest alcove.
\end{enumerate}
\end{proposition}
\leanref{src/Development/WeilBound/CoreStatements/AutomorphicKernels/MoriwakiBigness/MoriwakiBigTheorem/Outright.lean:9}{PositivityTheory.Pf1.HilbertSamuelMoriwaki.moriwakiBig\_outright}
\leanref{src/Development/WeilBound/CoreStatements/AutomorphicKernels/MoriwakiBigness/MoriwakiBigTheorem/Outright.lean:13}{PositivityTheory.Pf1.HilbertSamuelMoriwaki.hsUse\_outright}

\begin{proof}
Every single item in Clusters~I--VI is a published theorem verified on our smooth projective integral toroidal compactifications $\overline{\mathcal{S}}_K$ of PEL unitary Shimura varieties of signature $(p, q)$ over $\OO_F[1/N]$ (with neat level $K$, guaranteed by $2fk < \min(g_0, g_1) - 1$ as in \cref{thm:12:size-inputs}). Note in particular that the dead-code \texttt{ScholzeWeinsteinB} comparison is never invoked: Cluster~V supplies Fargues exactness for the Frobenius prime-jump axiom (T9) below and crystalline--de~Rham twisting in \cref{ch:14}, while flat determinant degrees in \cref{ch:14} use only Cluster~I (\texttt{ResK} + Serre GAGA)!
\end{proof}

\section{Flag Bundles, Well-Separated Weights, and the $p \ge 2q + 1$ Isotropic Bootstrap}
\label{sec:13:flag-bootstrap}

Let $S = \overline{\mathcal{S}}_K$ be an arithmetic toroidal compactification of a PEL unitary Shimura variety of signature $(p, q)$ ($q \ge 2$, $p \ge 2q + 1$) at a chosen real place $v_0 \mid \infty$ of $E^+$ (and definite signature $(p+q, 0)$ at all other real places of $E^+$). Let $\pi : \mathcal{F} \to S$ be the relative complete flag bundle of the rank-$q$ (or rank-$p$ and rank-$q$) Hodge vector bundles $\mathcal{E}^+ \oplus \mathcal{E}^-$, with relative flag dimension $d_{\mathrm{fl}}$.

To any dominant integral weight $u$ of $\UU(p) \times \UU(q)$ we associate:
\begin{itemize}
\item the tautological relative line bundle $\mathcal{L}(u)$ on $\mathcal{F}$, whose pushforward $\pi_* \mathcal{L}(u)$ is the automorphic vector bundle $\mathcal{V}_u$ on $S$ (by Borel--Weil--Bott);
\item its Weyl flag degree $\deg_{\mathrm{fl}}(u) := \int_{\mathcal{F}_s} c_1(\mathcal{L}(u))^{d_{\mathrm{fl}}} > 0$ on a geometric fiber $\mathcal{F}_s$; and
\item its unsquared Arakelov mean height at level cutoff $B$ on component $c$:
  \begin{equation}
  \label{eq:13:mean-height}
    \widehat{\mu}_{\mathcal{T}}(B, c; u) \;:=\; \frac{\mathcal{P}_c\bigl(B,\, \overline{\mathcal{L}}(u)^{d_{\mathrm{fl}} + 1}\bigr)}{(d_{\mathrm{fl}} + 1)\,\vol(c)\,\deg_{\mathrm{fl}}(u)},
  \end{equation}
  where $\mathcal{P}_c(B, \cdot)$ denotes the arithmetic intersection pairing on the level-$B$ root tower (\cref{def:13:root-tower} below).
\end{itemize}

A weight $u$ satisfies the \emph{well-separated weight hypothesis} (\texttt{WeightHyp}) if its consecutive weight gaps exceed the explicit Weyl-chamber threshold needed to make $\mathcal{L}(u)$ vertically ample and keep the boundary cusp curvature strictly subordinate to the interior Petersson metric. The exact numerical reason for the signature threshold $p \ge 2q + 1$ appears cleanly in the \emph{isotropic counting bootstrap}:

\begin{lemma}[Isotropic counting bootstrap ($p \ge 2q + 1$)]
\label{lem:13:isotropic-bootstrap}
Let $V$ be a Hermitian $E$-space of signature $(p, q)$ at $v_0$ ($q \ge 2$) and definite at all other real places, with $f := [E^+ : \QQ]$. On the toroidal boundary of the unitary Shimura variety $\overline{\mathcal{S}}_K$, the maximal $1$-dimensional boundary cusp strata carry an archimedean cusp-weight exponent $a := 2 f q$ (counting the real-rank-$q$ unipotent radical directions), while the interior Siegel--Koecher boundary-volume exponent is $\kappa := 2 f (p + q - 1)$. Consequently, the strict $L^1$ cusp-dominance inequality
\begin{equation}
\label{eq:13:bootstrap-ineq}
  \kappa \;\ge\; 3 a \qquad\Longleftrightarrow\qquad 2 f (p + q - 1) \;\ge\; 6 f q \qquad\Longleftrightarrow\qquad p \;\ge\; 2q + 1
\end{equation}
holds unconditionally whenever $p \ge 2q + 1$, guaranteeing that the cusp-boundary $L^1$ integral of the singular Chern--Weil curvature converges absolutely and is bounded by the uniform base anchor $C_{\mathrm{base}} = \max(C_h + 2, 0)$.
\end{lemma}
\leanref{src/Development/AutomorphicSpectrum/Definitions/DivisibilityTowers/CompactGapsGE.lean:22}{GlobalTheory.CompanionTheory.Ar.Top.degFl}
\leanref{src/Development/AutomorphicSpectrum/Definitions/DivisibilityTowers/CompactGapsGE.lean:73}{GlobalTheory.CompanionTheory.Ar.Top.WeightHyp}

\begin{proof}
In the toroidal compactification of $\UU(p, q)$~\cite{Lan2013}, a boundary component of Witt rank $r$ ($1 \le r \le q$) is fibered over a Shimura variety for $\UU(p - r, q - r)$ by an abelian scheme of relative dimension $r(p + q - 2r)$ and a torus of rank $r^2$. The invariant Petersson volume form near a cusp of rank $r$ behaves like $y_1^{-2f(p+q-1)} \cdots y_r^{-2f(p+q-2r+1)} dy_1 \cdots dy_r$ in Siegel-domain radial coordinates $y_i \to \infty$, whereas the singular log-canonical metric of a weight-$u$ automorphic bundle grows at most like $(\prod_{i=1}^r y_i)^{2fq}$ relative to the Hodge determinant. The integral of a test power $y_i^{3a - \kappa - 1}$ against $dy_i$ near $y_i = \infty$ converges with explicit constant $C_{\mathrm{base}}$ precisely when $\kappa \ge 3a$, which rearranges identically to $p + q - 1 \ge 3q \iff p \ge 2q + 1$!
\end{proof}

\section{The Four-Phase Abstract Flag-Bundle Tower Engine}
\label{sec:13:four-phase}

To eliminate the $O(\log B / B)$ error inherent in finite-level arithmetic intersection estimates and lock the exact proportionality of $\widehat{\mu}_{\mathcal{T}}(B, c; u)$ to the Weyl linear form $A_0 s_\sigma + b_0 t_\sigma$, we pass to an infinite \emph{arithmetic flag-bundle root tower} indexed by primes $B \to \infty$ (splitting in $E$ and adding root-level coverings of the universal abelian variety).

\begin{definition}[Arithmetic flag-bundle root tower and its three axiom clusters]
\label{def:13:root-tower}
An \emph{arithmetic flag-bundle root tower} $\mathcal{T} = \{(\mathcal{F}_B, \overline{\mathcal{L}}_B(u))\}_{B}$ over the unitary Shimura base with Weyl symmetry group $\Gamma$ and level-volume scaling $\vol(c)$ satisfies three compact clusters of axioms (consolidating the thirteen tower axioms \texttt{(T1)}--\texttt{(T13)} of the formal archive):
\begin{enumerate}[label=\textup{\textbf{Ax.~\arabic*.}}]
\item \emph{(Analytic and vertical nefness package: \texttt{(T2)}--\texttt{(T7)}.)} On every geometric fiber of $\mathcal{F}_B$, the root line bundles are vertically nef; the top arithmetic intersection $\mathcal{P}_c(B, \cdot)$ is symmetric, multilinear, and continuous in the weight vector $u$, satisfying Zhang's successive-minima / Moriwaki bigness inequality (\cref{prop:13:six-clusters}, Cluster~II) with uniform defect constant $C$.
\item \emph{(Weyl symmetry, Frobenius prime jumps, and equal component volumes: \texttt{(T1)}, \texttt{(T8)}--\texttt{(T10)}, \texttt{(T12)}.)} The Weyl group $\Gamma$ acts transitively on chambers with antisymmetric wall reflection; passing from level $B$ to a prime-power cover $p B$ shifts the root line bundle by the exact \emph{Fargues Frobenius prime-jump formula} (T9) (\cref{prop:13:six-clusters}, Cluster~V); and all connected components $c$ share equal geometric volume (T12).
\item \emph{(Two quantitative anchors: \texttt{(T11)} and \texttt{(T13)}.)} On the flat line (vanishing curvature along the diagonal), the arithmetic degree is bounded by $C' = 4$ times the height scale $\Lambda$ (T11); on the base cusp test weight $P_{\mathrm{test}}(j)$, \cref{lem:13:isotropic-bootstrap} bounds the boundary contribution by $C_{\mathrm{base}} = \max(C_h + 2, 0)$ (T13).
\end{enumerate}
\end{definition}
\leanref{src/Development/AutomorphicSpectrum/Definitions/DivisibilityTowers/Proportional.lean:10}{GlobalTheory.CompanionTheory.Ar.Top.RootTower}
\leanref{src/Development/AutomorphicSpectrum/Proofs/GeometryProofs/ExistsGoodRootTower.lean:77}{GlobalTheory.CompanionTheory.Ar.TopTower.exists\_good\_tower}
\leanref{src/Development/AutomorphicSpectrum/Proofs/GeometryProofs/RootLinesWeightBoundOfFree.lean:35}{GlobalTheory.CompanionTheory.Ar.HT.rootModel\_nefAt\_of\_hodgeCpt}
\leanref{src/Development/WeilBound/CoreStatements/FlagImmersions/GlobalModelTowerFifthCondition.lean:22}{RootFieldNumerics.global\_model\_tower\_fifth\_condition}

We now prove the abstract flag-bundle tower theorem in \emph{four sharp phases} (distilling the eight-step archive proof while preserving every exact constant!).

\begin{proposition}[Four-phase abstract flag-bundle tower mean-height engine]
\label{prop:13:tower-engine}
Let $\mathcal{T}$ be an arithmetic flag-bundle root tower satisfying \textbf{Ax.~1--Ax.~3} of \cref{def:13:root-tower} with parameters $(C, C' = 4, C_u = 4, C_{\mathrm{base}})$, and let $K_{\mathrm{tot}} := 2C + 3C' + 5$. Then for every well-separated dominant weight $u$ and every connected component $c$ at level $B \ge B_0$:
\begin{equation}
\label{eq:13:tower-engine}
  \widehat{\mu}_{\mathcal{T}}(B, c; u)
  \;\ge\; -\bigl(K_{\mathrm{tot}} + C_{\mathrm{base}}\bigr)\,\Lambda\,\|u\|_1 \;-\; \frac{C_u \log B}{B}\,\|u\|_1.
\end{equation}
\end{proposition}
\leanref{src/Development/AutomorphicSpectrum/Proofs/DivisibilityProofs/TowerWeilTheorem.lean:120}{GlobalTheory.CompanionTheory.Ar.Top.thmW\_tower}

\begin{proof}
We execute the proof in four crisp phases, tracking every explicit constant:

\emph{Phase~1: Wall vanishing and Weyl rigidity.}
On any wall of the Weyl chamber (where two adjacent entries of the weight vector $u$ coincide), apply Zhang's essential-minimum squeeze (\cref{prop:13:six-clusters}, Cluster~II) with perturbation parameter $\theta := 5 f / B$ and step $s := \theta^2$. Because $\deg_{\mathrm{fl}}(u)$ vanishes to order at least $1$ on the wall while the perturbed arithmetic intersection grows linearly in $\theta$, taking $s = \theta^2$ forces the wall residual of $\mathcal{P}_c(B, \cdot)$ to vanish up to $O(\log B / B)$! Averaging over the transitive Weyl group $\Gamma$ via antisymmetry across the walls (Ax.~2) forces the leading $B$-asymptotic of $\widehat{\mu}_{\mathcal{T}}(B, c; u)$ to be a linear combination $A_0 s_\sigma(u) + b_0 t_\sigma(u)$ of the two $\Gamma$-invariant basis forms (the interior trace $s_\sigma$ and the flat-determinant shift $t_\sigma$).

\emph{Phase~2: Exact prime-cutoff identity via Frobenius jumps ($B \to \infty$).}
Evaluate the Fargues prime-jump identity (T9) at two distinct primes $p_1 < p_2$ dividing large levels $B \to \infty$. This yields a $2 \times 2$ triangular linear system on the pair $(p_1 Q_1, p_2 Q_2)$ of level-jump quotients; letting $B \to \infty$ kills the $O(\log B / B)$ wall error \emph{identically} and expresses the exact coefficients $(A_0, b_0)$ in terms of the flat-line anchor (T11). Specializing to $k = 1$ with $C' = 4$ pins the exact linear combination
\begin{equation}
\label{eq:13:jA-Nb}
  \bigl|j\, A_0 + N\, b_0\bigr| \;\le\; 4 N\,\Lambda.
\end{equation}

\emph{Phase~3: Base anchoring via the $p \ge 2q + 1$ isotropic bootstrap.}
Evaluate the linear form on the canonical boundary test weight $P_{\mathrm{test}}(j)$, chosen so that the Shimura shift sum $\sum \mathrm{Sh}(P_{\mathrm{test}}(j)) = 0$ vanishes identically. By \cref{lem:13:isotropic-bootstrap} ($p \ge 2q + 1$, Ax.~3), the boundary cusp contribution is bounded below by $-C_{\mathrm{base}}\Lambda$, while Moriwaki bigness (Ax.~1, constant $C$) and \eqref{eq:13:jA-Nb} ($C' = 4$) combine to anchor the interior slope coefficient:
\[
  A_0 \;\ge\; -\bigl(2(C + C') + 3 + C_{\mathrm{base}}\bigr)\,\Lambda.
\]

\emph{Phase~4: Dominant-weight rearrangement and equal component volumes.}
Rewrite the Weyl linear form identically as
\[
  A_0\, s_\sigma(u) + b_0\, t_\sigma(u)
  \;=\; A_0\!\left(s_\sigma(u) - \frac{j}{N}\, t_\sigma(u)\right) \;+\; \bigl(j A_0 + N b_0\bigr)\,\frac{t_\sigma(u)}{N}.
\]
For every well-separated dominant weight $u$, $s_\sigma(u) - \frac{j}{N} t_\sigma(u) \ge 0$ and $|t_\sigma(u)/N| \le \|u\|_1$. Substituting the Phase~3 lower bound for $A_0$ and \eqref{eq:13:jA-Nb} for $|j A_0 + N b_0|$, dividing by the equal component volume $\vol(c)$ (T12), and restoring the finite-level residual $C_u \frac{\log B}{B}\|u\|_1$ yields \eqref{eq:13:tower-engine} with $K_{\mathrm{tot}} = 2C + 3C' + 5$.
\end{proof}

\section{PEL Instantiation: The Volume Theorem in Signature $(p, q)$}
\label{sec:13:volume-pq}

Instantiating \cref{prop:13:tower-engine} on our integral PEL unitary Shimura flag tower (constructed using \cref{prop:13:six-clusters}, Clusters~I--VI, with vertical nefness of root line bundles on every geometric fiber and global point-height bounds along the tower) establishes the flagship theorem of \cref{ch:13}.

\begin{theorem}[Arakelov Mean Height Lower Bound in Signature $(p, q)$ ($p \ge 2q + 1$)]
\label{thm:13:volume-pq}
Let $E/E^+$ be our host CM extension, let $V$ be a Hermitian $E$-space of signature $(p, q)$ at a distinguished real place $v_0 \mid \infty$ of $E^+$ (with $q \ge 2$ and $p \ge 2q + 1$) and definite signature $(p+q, 0)$ at all other real places, and let $\overline{\mathcal{S}}_K$ be a smooth projective integral toroidal compactification of the associated PEL unitary Shimura variety at neat level $K$. Then there exists an admissible arithmetic flag-bundle root tower $\mathcal{T}$ over $\overline{\mathcal{S}}_K$ (\cref{def:13:root-tower}) and an absolute constant $C_{\mathrm{ht}} > 0$ such that for every well-separated dominant weight $u$ and every connected component $Z \in \mathrm{Comp}(\overline{\mathcal{S}}_K)$, the unsquared Arakelov mean height obeys
\begin{equation}
\label{eq:13:vol-pq}
  \widehat{\mu}_{\mathcal{T}}(\lfloor B \rfloor, \kappa(Z); u)
  \;\ge\; -C_{\mathrm{ht}}\,\mathcal{E}_{\mathrm{geo}}\,\|u\|_1 \;+\; 2\operatorname{cost}_2(l, B)\|u\|_1,
\end{equation}
where $\mathcal{E}_{\mathrm{geo}}$ is the explicit arithmetic base-height envelope (proportional to $\log\rd_E + \log(k b)$) and $\operatorname{cost}_2(l, B)$ is the dyadic bad-prime-$2$ metric adjustment (still in unsquared normalization, ready for factor-of-$2$ conversion in \cref{ch:14}).
\end{theorem}
\leanref{src/Development/AutomorphicSpectrum/Statements/TowerExistence/GoodTowerExistsX.lean:28}{GlobalTheory.CompanionTheory.Ar.Top.VolumeTheoremX}

\begin{proof}
Construct the integral root flag-bundle tower $\mathcal{T}$ over $\overline{\mathcal{S}}_K$ using Cluster~IV (Lan's integral toroidal compactifications~\cite{Lan2013} and Mumford's good singular Chern--Weil metrics~\cite{Mumford1977}) and Cluster~V (Fargues's exact Frobenius jump formula~\cite{Fargues2010} for the prime-power $p$-isogeny covers of the universal PEL abelian scheme). Vertical nefness of the root line bundles on every geometric fiber follows from Borel--Weil--Bott and the Hodge-compactness positivity of Cluster~III; the global point-height fifth condition follows from Cluster~II (cycle heights and Chow forms~\cite{BostGilletSoule1994,Zhang1995}) and Cluster~III (Ohsawa--Takegoshi $L^2$-extension~\cite{OhsawaTakegoshi1987} and Angehrn--Siu effective point separation~\cite{AngehrnSiu1995}). Since $p \ge 2q + 1$, \cref{lem:13:isotropic-bootstrap} discharges axiom (T13). Applying \cref{prop:13:tower-engine} at cutoff $B$ large enough that $(C_u \log B)/B \le \mathcal{E}_{\mathrm{geo}}$ and separating the dyadic prime-$2$ metric adjustment $\operatorname{cost}_2(l, B)$ yields \eqref{eq:13:vol-pq} unconditionally.
\end{proof}
